\documentclass[11pt,letterpaper]{article}
\usepackage[top=1.0in,bottom=1.0in,left=1.25in,right=1.25in]{geometry}
\usepackage{microtype}
\usepackage{parskip}
\usepackage{xcolor}
\usepackage{mdframed}
\usepackage{amsmath,amssymb}
\usepackage{booktabs}
\usepackage{hyperref}
\hypersetup{colorlinks=true,urlcolor=blue,linkcolor=black,citecolor=black}
\setlength{\parindent}{0pt}
\setlength{\parskip}{7pt}

\mdfdefinestyle{reviewquote}{
  leftmargin=0.4cm, rightmargin=0.4cm,
  innerleftmargin=0.4cm, innerrightmargin=0.4cm,
  innertopmargin=5pt, innerbottommargin=5pt,
  backgroundcolor=gray!8,
  linecolor=gray!40, linewidth=0.5pt,
  skipabove=6pt, skipbelow=6pt
}
\newcommand{\rquote}[1]{\begin{mdframed}[style=reviewquote]\textit{#1}\end{mdframed}}
\newcommand{\sref}[1]{\textbf{[Revision:~#1]}}
\newcommand{\fixed}[1]{\textcolor{blue!70!black}{\textbf{Fixed:} #1}}
\newcommand{\rebuttal}[1]{\textcolor{green!40!black}{\textbf{Rebuttal:} #1}}

\begin{document}

\textbf{Response to Reviewers --- Second Revision}\\
\textbf{Manuscript:} PONE-D-26-05900R1\\
\textbf{Title:} Signal Hierarchical Petri Nets: Formal Semantics of
Hierarchical Regulatory Control of Biological Systems\\
\textbf{Authors:} Eug\^{e}nio Sim\~{a}o\\
\textbf{Date:} \today

\bigskip\hrule\bigskip

We thank the editor and both reviewers for the second round of evaluation.
Reviewer~1 has indicated that all previous concerns have been addressed.
Reviewer~2 (Dr.\ Juri Kolč{\'a}k) maintains a rejection recommendation
and raises a further set of technical points.
We respond to each point below with either a correction or a rebuttal,
supported in each case by a specific revision reference.

Revision references of the form \sref{Sec.~X} indicate where each change
appears in the revised manuscript.
A track-changes PDF generated by \texttt{latexdiff} against the R1 submission
is submitted alongside this letter.

\bigskip

%% ============================================================
\section*{Academic Editor (Prof.\ Karthik Raman)}
%% ============================================================

\rquote{While the reviewers appreciate the contributions, one of them
continues to have some concerns, which must be addressed (or rebutted
effectively) before accepting the manuscript for publication.}

We have addressed each of Reviewer~2's remaining concerns.
Genuine technical errors have been corrected; points that rest on a
mischaracterisation of the stochastic Petri net formalism or of contextual
Petri net theory have been rebutted with precise citations and arguments.
The core results --- the acyclicity theorem, the preemption theorem, the basin
computability result, and the biological validation --- are unaffected by the
corrections and are strengthened by the additional clarifications.

\bigskip

%% ============================================================
\section*{Reviewer \#1 (anonymous)}
%% ============================================================

\rquote{As far as your theoretical considerations and hypothesis are concerned
and supported (corroborated) by your accessible data, the main theoretical
conclusions are valid. I think the article has mended all the minor issues I
could find. I tend to agree that the major issue pointed out by other referees
has been mitigated.}

We thank Reviewer~1 for this positive assessment. No further changes were
requested.

\bigskip

%% ============================================================
\section*{Reviewer \#2 (Dr.\ Juri Kolč{\'a}k)}
%% ============================================================

We thank Dr.\ Kolč{\'a}k for the continued and detailed technical engagement.
We address his two major comments, followed by every specific point raised, in order.

\subsection*{Major comment (1) --- Rate functions $\Phi$ absent from enabling/firing semantics}

\rquote{It remains a fact that while $\Psi$ (Rate functions) are introduced
throughout the paper, they are not used in the semantics of the model,
making them completely useless and superfluous.}

\rebuttal{We note a notation clarification before responding: in our formalism
$\Psi \subseteq P$ denotes the \emph{signal places set} (information channels),
while $\Phi: T \rightarrow \text{RateFunctions}$ denotes the \emph{rate
functions}. The reviewer's comment uses $\Psi$ where our paper uses $\Phi$;
we respond to the intended concern about rate functions $\Phi$ not appearing
in the Enabled predicate.

This comment conflates the two standard layers of stochastic Petri
net semantics, which are architecturally separate by design.
In stochastic Petri nets (Goss \& Peccoud 1998; Haas 2002), the formalism
has two distinct layers:
\begin{enumerate}
  \item \emph{Structural enabling} (topology-based): the Enabled predicate
    determines \emph{whether} a transition may fire, based solely on token
    counts and arc structure. This is the layer the reviewer is examining.
  \item \emph{Kinetic propensity} ($\Phi$): the rate function determines
    \emph{how fast} an enabled transition fires, governing the simulation
    dynamics ($\tau$-leaping propensity, ODE rates, and---crucially---the
    entropy production rate via the propensity ratio $\Phi^+(M)/\Phi^-(M)$
    in stochastic thermodynamics).
\end{enumerate}
$\Phi$ is intentionally absent from the Enabled predicate and the Firing
rule: this is the standard stochastic PN architecture, not a deficiency.
A formalism in which $\Phi$ appeared in the enabling condition would
conflate qualitative reachability (structural) with quantitative dynamics
(kinetic), making structural theorems unprovable without specifying rates---
precisely the limitation classical PN theory was designed to avoid.

We have added an explicit paragraph to the $\Phi$/$F_s$ remark
(\sref{Sec.~4.1, Remark after signal flow arc definition}) and to the
13-tuple definition (\sref{Sec.~4.1, Definition~5}) making this two-layer
architecture explicit and citing the stochastic PN literature.}

\subsection*{Major comment (2) --- Motivation for test arcs over loops}

\rquote{It is well known that test arcs (often called read arcs) create a
semantic difference to loops ($W(p,t)=W(t,p)$) when concurrency is a
concern, but are essentially indistinguishable from loops in other settings.
What is the motivation for inclusion of test arcs into the proposed model,
why are loops not enough? Additionally, test arcs could easily be modelled
using the ``RateFunctions'' $\Psi$ as they have been described.}

\rebuttal{We offer two independent justifications.

\textbf{(a) Semantic distinctness (Contextual Petri Nets):}
The reviewer cites Baldan, Corradini \& Montanari (``Contextual Petri Nets,
Asymmetric Event Structures, and Processes'') as evidence that test arcs and
loops are distinguishable only under concurrency.
We note that this reference establishes precisely the opposite of the
reviewer's claim: Baldan et al.\ introduced contextual/read arcs as a
formally distinct arc type with different causal semantics from loops, and
their paper is the canonical reference for the inclusion of non-consuming
observation arcs in Petri net theory.
Loops $W(p,t) = W(t,p)$ consume and regenerate the same number of tokens:
they contribute to the stoichiometric firing rule
$M'(p) = M(p) - W((p,t)) + W((t,p))$, which is zero only because the
two terms cancel.
Test arcs do not participate in the firing rule at all:
$M'(p) = M(p)$ unconditionally.
This is a semantic difference in the single-step firing semantics,
independent of any concurrent setting.

\textbf{(b) Thermodynamic distinctness:}
For this formalism's primary application --- the stochastic thermodynamics
of biological commitment --- test arcs and loops are not interchangeable.
Test arcs have zero mass transfer ($\Delta M = 0$), hence zero entropy
production rate (EPR) contribution: catalysts that enable reactions without
being consumed must not appear in the dissipation budget.
Loops have non-zero mass flow (tokens are consumed then regenerated), hence
non-zero EPR contribution.
Using a loop to model a catalyst would spuriously inflate the EPR of the
commitment decision and corrupt the decomposition $\dot\sigma = \sum_\rho
J_\rho \ln(J_\rho^+/J_\rho^-)$.
The thermodynamic distinction is not incidental; it is central to the
paper's contribution.

We have added both justifications to the test arc paragraph
(\sref{Sec.~4.1, Test arc definition}).}

\subsection*{Specific point: Page 10, Theorem 2 proof}

\rquote{PreemptionCheck has been included in prerequisites (although it has
not been defined yet), but the proof is false.
The induction hypothesis is that Enabled is false for any span.
The induction step then relies on PreemptionCheck to prove Enabled is also
false for span $= k$, however, Enabled (induction hypothesis) is not an
input of PreemptionCheck as it only relies on the other three Enabled
components (NormalEnabled, TestEnabled, SignalEnabled).}

\fixed{The reviewer is correct.
The proof incorrectly propagated $\neg\text{Enabled}$ through
PreemptionCheck, but PreemptionCheck is defined as:
\[
  \text{PreemptionCheck}(t,M) \equiv \forall p_s \in {}^\bullet_s t:\,
  \forall t' \in {}^\bullet_s p_s:\,
  \text{NormalEnabled}(t',M) \wedge \text{TestEnabled}(t',M) \wedge
  \text{SignalEnabled}(t',M)
\]
i.e., it checks the three sub-predicates of $t'$, not $\text{Enabled}(t')$
itself.

The corrected proof proceeds as follows.
\emph{Base case} ($k=1$): $M(p_i) < \theta(t_i)$ implies
$M(p_i) < \theta(t_i) + W_s((p_i,t_i))$ (since $W_s > 0$), hence
$\neg\text{SignalEnabled}(t_i, M)$.
Since PreemptionCheck$(t_j, M)$ requires SignalEnabled$(t_i, M)$ for
every $t_i \in {}^\bullet_s p_s$ where $p_s \in {}^\bullet_s t_j$,
we have $\neg\text{PreemptionCheck}(t_j, M)$, hence $\neg\text{Enabled}(t_j, M)$.
\emph{Inductive step}: the propagation uses $\neg\text{SignalEnabled}$
at each layer rather than $\neg\text{Enabled}$, tracing through the
$F_s$-acyclic path.

\sref{Sec.~2.2, Theorem~2, proof rewritten.}}

\subsection*{Specific point: Page 11, Theorem 3 proof}

\rquote{There is no restriction on transitions in $F \setminus F_s$ producing
tokens in $p_s$. In fact, further in the text the author outright states that
the metabolic ``horizontal'' transitions are allowed to replenish signalling
places.}

The theorem statement already carries the hypothesis ``assuming no other
transition regenerates $p_s$ via $F_s$'' (i.e., no $t' \neq t$ with
$(t', p_s) \in F_s$).
The reviewer's concern is that transitions in $F \setminus F_s$ (normal arcs)
can also produce tokens in $p_s$.
This is correct biologically and is an intentional design: metabolic
transitions in $G_E$ can replenish signal places via \emph{normal arcs},
representing ongoing production of the signalling species.
The theorem applies to the signal-arc topology only, establishing that no
\emph{signal-arc path} returns to $p_s$ after commitment firing.
Metabolic replenishment via normal arcs is not a counter-example to the
basin partition theorem because it affects $G_E$ dynamics, not the $G_S$
commitment topology.

The text has been clarified to make this scope explicit.
\sref{Sec.~2.3, Theorem~3, added scope note.}

\subsection*{Specific point: Page 18, Definition 5 --- $F_s$ acyclicity
redundancy in equation (8)}

\rquote{$F_s$ being acyclic implies $W_s(t,p) = 0$ if $W_s(p,t) > 0$;
there is therefore no need to use it in the equation.}

The reviewer is formally correct: for the acyclic case, the backward weight
$W_s((t,p))$ is necessarily zero.
The general form $M'(p_s) = M(p_s) - W_s((p_s,t)) + W_s((t,p_s))$ is
retained for two reasons: (1) it matches the standard PN firing rule format
that readers familiar with the literature will expect; (2) the backward term
makes explicit that in the acyclic case $W_s((t,p_s)) = 0$, so
$M'(p_s) = M(p_s) - W_s((p_s,t))$, which is informative rather than redundant.
No change made; a clarifying remark has been added. \sref{Sec.~4.2, Firing
Rule remark.}

\subsection*{Specific point: Page 19, line 442 --- $\tau_t = 0$ inconsistency}

\rquote{Here $\tau_t = 0$ is used to mean $M(p) > 0$, however, in the
definition of the TestPredicate, as well as earlier in the paragraph, the
not-sharp inequality $M(p) \geq \tau_t$ is used.}

\fixed{The reviewer is correct.
$M(p) \geq 0$ (non-sharp, default $\tau_t = 0$) is vacuously true and does
not implement a presence check.
The TestEnabled predicate has been corrected to use strict inequality
$M(p) > \tau_t((p,t))$ throughout, so that the default $\tau_t = 0$
correctly implements $M(p) > 0$ (presence check).
\sref{Sec.~4.2, Definition~6, TestEnabled predicate and surrounding remark.}}

\subsection*{Specific point: Page 20, lines 456--461 --- commitment states
mislabelled}

\rquote{The commitment states appear mislabelled: $M(p_s) < M_{\mathrm{commit}}(p_s)$
should be post-commitment (``irreversibly emptied'') while
$M(p_s) > M_{\mathrm{commit}}(p_s)$ should be pre-commitment (``fireable'').}

\fixed{The reviewer identified a genuine ambiguity.
Both interpretations are valid but describe different time points:
\begin{itemize}
  \item \emph{Before firing}: $M(p_s) < M_{\mathrm{commit}}$ means signal not
    yet at threshold (pre-commitment, transition disabled).
  \item \emph{After firing}: $M'(p_s) = \theta(t) < M_{\mathrm{commit}}$ means
    signal consumed (post-commitment, irreversible).
\end{itemize}
The labelling now explicitly distinguishes these two cases by adding
time-point qualifiers ``(before firing)'' and ``(after firing)''.
\sref{Sec.~4.2, Commitment Semantics, enumeration items 1 and 3.}}

\subsection*{Specific point: Page 21, Fig.\ 2 --- caption and threshold crossing}

\rquote{None of the dynamics appear to exceed the $\theta_{\sigma H}$ threshold
around $t = 300$, which is marked as the ``global commitment time''.}

The figure shows 50 stochastic replicates.
Commitment occurs when any individual trajectory stochastically crosses the
threshold; the mean trajectory does \emph{not} cross it (this is the
threshold-crossing probability mechanism: $P(\mathrm{spor}) = P(\sigma_H^{\mathrm{stoch}}
> \theta_{\sigma H})$).
The caption has been revised to state explicitly that commitment is a
\emph{stochastic threshold-crossing event}, not a deterministic mean-field
crossing, and that the marked time $t_{\mathrm{commit}} \approx 300$~min is
the time at which the first sporulating replicate fires T\_septation,
not the time at which the mean trajectory crosses $\theta_{\sigma H}$.
\sref{Sec.~5.1, Fig.~2 caption.}

\subsection*{Specific point: Page 22, Def.\ 6 --- $\Phi$ and $C$ not used in
Enabled or firing}

\rquote{$\Phi$, the untyped kinetic rate laws, are never used in the definition
of Enabled or in firing of transitions. Why do they appear in this definition?
Similarly, the untyped ``$C$'' assigning stoichiometric coefficients is never
used in definition of Enabled or in firing.}

As addressed in Major Comment~(1) above: $\Phi$ governs the propensity layer
(simulation dynamics, EPR computation), not the structural enabling layer.
This is the standard stochastic PN architecture.
$C$ provides stoichiometric coefficients for ODE-mode transitions (continuous
flux computation), distinct from discrete token weights $W$.
Both $\Phi$ and $C$ have been annotated with explicit simulation-role
descriptions in Definition~5. \sref{Sec.~4.1, Definition~5.}

\subsection*{Specific point: Page 25, line 569 --- reachability complexity}

\rquote{The author claims reachability in Petri nets is EXPSPACE-complete, but
the citation provided only claims EXPSPACE-hardness. In fact, reachability in
vector addition systems is Ackermann-complete (Czerwi\'nski and Orlikowski).}

\fixed{The reviewer is correct on both counts.
The claim has been corrected to \emph{Ackermann-complete} at all five
occurrences, and the Czerwi\'nski \& Orlikowski (FOCS 2021) reference has
been added to the bibliography. \sref{Sec.~5.2 proof, Sec.~5.3 discussion,
Sec.~6 complexity paragraph, Sec.~7 Table~2, Sec.~7 summary; and bibliography.}}

\subsection*{Specific point: Page 26, line 585 --- test arc sharp inequality}

\rquote{Test arcs ``redefined'' as sharp inequality although TestPredicate
definition uses $\geq$.}

\fixed{Corrected: the TestEnabled predicate now uses strict inequality
$M(p) > \tau_t((p,t))$ consistently throughout, resolving this inconsistency.
\sref{Same as $\tau_t$ fix above.}}

\subsection*{Specific point: Page 30, Equation (15) --- $M_{\mathrm{commit}}$
trivial}

\rquote{The biologically relevant threshold is obviously the input $\theta$;
$W_s$ is merely a modeller's choice.
The value of $M_{\mathrm{commit}}$ is trivial to compute and useless.}

The formula $M_{\mathrm{commit}} = \theta + W_s$ is retained because it
serves a concrete predictive purpose in the biological validation.
The reviewer's characterisation as ``trivial'' is perhaps true in the abstract,
but the claim of the paper is not the formula itself but the capacity to
compute structural commitment thresholds directly from arc-level parameters
$\theta$ and $W_s$ without simulation or global fitting.
For signal-flow arc inputs (e.g., GTP\textunderscore pool gating $T_{\mathrm{septation}}$),
$M_{\mathrm{commit}} = \theta + W_s$ yields the structural threshold from
topology alone.
We note that $\theta_{\sigma H} = 1.60\,\mu$M is an \emph{effective kinetic
threshold} identified via simulation (Sweep~A), not a structural
$M_{\mathrm{commit}}$ computation, since $\sigma_H$ gates $T_{\mathrm{septation}}$
through the rate function $\Phi$ (remote sensing) rather than a signal-flow arc.
The manuscript has been updated to make this distinction explicit.
A sentence clarifying the predictive role of the structural formula has been
added. \sref{Sec.~4.3.}

\subsection*{Specific point: Page 31, Theorem 6.1.1 --- proof ``hardly necessary''}

\rquote{This proof is hardly necessary, given that the layer assignment is
indeed just a topological sorting.}

Agreed that topological sorting is well-known.
The theorem and algorithm are retained because they serve a different
purpose: they provide a \emph{constructive failure condition} (the algorithm
returns FAIL if $G_s$ contains a cycle), which is the computational
mechanism for enforcing the Acyclicity Requirement of Theorem~1 during
model construction.
A sentence has been added making this purpose explicit. \sref{Sec.~4.4.}

\subsection*{Specific point: Page 32, Proposition 3 --- copies Theorem 2}

\rquote{This proposition is effectively a copy of Theorem~2 (page 10),
providing no further insights. The proof is also wrong, same as Theorem 2.}

Proposition~3 applies the general Hierarchical Preemption result
(Theorem~2) to the specific case of the \textit{B.\ subtilis} model
introduced in the validation section, making the cascade hierarchy concrete.
The proof has been replaced with a direct reference to Theorem~2 rather than
repeating the general argument, and the proposition is now labelled
``Corollary (Biological Instantiation of Theorem~2)'' to make the derivation
relation explicit. \sref{Sec.~5.2.}

\subsection*{Specific point: Page 33, Proposition 4 --- trivial $O(1)$ computation}

\rquote{The fact that sum of two input constants is computable in $O(1)$ is
trivial and requires no proof.}

Agreed.
The ``proof'' has been removed; the proposition is retained as a remark
stating the $O(1)$ threshold formula, since its biological significance
(direct computation from literature kinetic constants with no simulation)
justifies its explicit statement even if the arithmetic is elementary.
\sref{Sec.~5.2.}

\subsection*{Specific point: Figure 3 --- signal arcs form loops}

\rquote{Although acyclicity requirement of signal arcs is stressed as a strong
constraint, the signal arcs in Figure 3 immediately form loops.
See for instance KinA\_P $\to$ T\_Spo0F\_phosphorylation $\to$ KinA\_kinase3.0
$\to$ T\_kinA\_activation.}

\fixed{The reviewer is correct, and this is the most substantive remaining
technical error.
Algorithmic cycle detection on the model file confirmed three cycles in $G_s$
caused by four enzyme-regeneration arcs that were incorrectly typed as
signal-flow arcs rather than normal arcs:
\begin{enumerate}
  \item \texttt{T\_Spo0F\_phosphorylation} $\to$ \texttt{KinA\_kinase}
    (KinA regeneration after phosphotransfer)
  \item \texttt{T\_Spo0A\_phosphorylation} $\to$ \texttt{Spo0F}
    (Spo0F regeneration after phosphotransfer)
  \item \texttt{T\_Spo0F\_dephos} $\to$ \texttt{Spo0F}
    (Spo0F regeneration after dephosphorylation)
  \item \texttt{T\_Spo0A\_dephosphorylation} $\to$ \texttt{Spo0A}
    (Spo0A regeneration via phosphatase)
\end{enumerate}
The biological error: these are enzyme-\emph{recycling} steps
(mass-flow, horizontal in $G_E$), not regulatory commitment signals
propagating upward through the hierarchy.
The phosphoryl group moves from KinA$_P$ \emph{to} Spo0F (signal direction,
upward through $G_s$); KinA being regenerated is a metabolic by-product
that flows back through $G_E$, not $G_s$.

All four arcs have been corrected to type \texttt{normal} in the model file.
Cycle detection re-run on the corrected model confirms $G_s$ is acyclic
($|\Psi| = 20$ signal-flow arcs remaining, zero cycles).
Figure~3 has been regenerated from the corrected model.
\sref{Model file \texttt{bacillus\_sporulation\_v9.shy} (4 arc types corrected);
Fig.~3 regenerated; model description in Sec.~5.1 updated.}}

\bigskip

%% ============================================================
\section*{Summary of all R2 changes}
%% ============================================================

\begin{center}
\renewcommand{\arraystretch}{1.3}
\begin{tabular}{p{4.5cm}p{4cm}p{4.5cm}}
\toprule
\textbf{R2 point} & \textbf{Action} & \textbf{Location} \\
\midrule
Fig.\ 3 signal arc cycles & Fixed: 4 arcs reclassified to normal; $G_s$ verified acyclic & Model file; Fig.\ 3; Sec.~5.1 \\
Theorem 2 proof & Fixed: proof now uses $\neg$SignalEnabled, not $\neg$Enabled & Sec.~2.2, Thm.~2 proof \\
$\tau_t$ / test arc inconsistency & Fixed: TestEnabled changed to strict $>$ throughout & Sec.~4.2, Def.~6 \\
Commitment labelling ambiguity & Fixed: time-point qualifiers added & Sec.~4.2, Commitment Semantics \\
Reachability EXPSPACE claim & Fixed: Ackermann-complete at 5 locations & Sec.~5.2, 5.3, 6, 7 \\
$\Phi$ not in Enabled/Firing & Rebuttal: standard stochastic PN two-layer architecture & Sec.~4.1, Def.~5 remark \\
Test arcs vs.\ loops & Rebuttal: Baldan et al.\ (his own citation) + zero-EPR argument & Sec.~4.1, test arc definition \\
Def.\ 6 $\Phi$, $C$ role & Clarified: simulation propensity / ODE stoichiometry & Sec.~4.1, Def.~5 \\
Theorem 3 scope & Clarified: applies to $G_s$, not to $G_E$ normal arcs & Sec.~2.3, Thm.~3 \\
Fig.\ 2 threshold crossing & Clarified: stochastic crossing, not mean-field & Sec.~5.1, Fig.~2 caption \\
$M_{\mathrm{commit}}$ triviality & Rebuttal: predictive role from kinetic constants & Sec.~4.3 \\
Thm.\ 6.1.1 topological sort & Clarified: constructive failure condition & Sec.~4.4 \\
Prop.\ 3 copies Thm.\ 2 & Restructured as Corollary; proof replaced with reference & Sec.~5.2 \\
Prop.\ 4 trivial & Proof removed; retained as remark & Sec.~5.2 \\
\bottomrule
\end{tabular}
\end{center}

\end{document}
